\chapter{Исследовательская часть}

\section{Технические характеристики}

Технические характеристики устройства, на котором проводились замеры:

\begin{itemize}
	\item ARM Cortex-M7 core 216 МГц;
	\item Оперативная память: 2 Мб.
\end{itemize}

\section{Временные характеристики}

В таблице \ref{tbl:time} приведены временные характеристики, полученные в результате замеров времени алгоритмов. Для каждого размера квадратной матрицы замер проводился 100 раз, при этом для каждой итерации создавались 2 случайные матрицы заданного размера и использовались как аргументы к алгоритмам. Итоговое время взято как среднее арифметическое всех полученных замеров.

\begin{table}[H]
    \caption{Сравнение времени выполнения для различных методов умножения матриц}
    \label{tbl:time}
    \begin{tabularx}{\textwidth} { 
      | >{\centering\arraybackslash}X 
      | >{\centering\arraybackslash}X 
      | >{\centering\arraybackslash}X 
      | >{\centering\arraybackslash}X | }
        \hline
        \textbf{Размер матрицы, N$\times$N} & \textbf{Стандартный, нс} & \textbf{Виноград, нс} & \textbf{Виноград оптимизированный, нс} \\
    \end{tabularx}
    \begin{tabularx}{\textwidth} { 
      | >{\centering\arraybackslash}X 
      | >{\raggedleft\arraybackslash}X 
      | >{\raggedleft\arraybackslash}X 
      | >{\raggedleft\arraybackslash}X | }
        \hline
        15 & 2.50 & 2.71 & 2.54 \\
        \hline
        20 & 5.86 & 6.11 & 5.73 \\
        \hline
        25 & 11.33 & 11.74 & 11.12 \\
        \hline
        30 & 19.45 & 19.86 & 18.85 \\
        \hline
        35 & 31.01 & 31.26 & 29.93 \\
        \hline
    \end{tabularx}
\end{table}

Полученные замеры также можно увидеть на рисунке \ref{fig:time}:

\begin{figure}[H]
	\centering
	\includegraphics[width=1\textwidth]{time}
	\caption{График зависимости времени работы алгоритмов от размера матриц}
	\label{fig:time}
\end{figure}

Как видно из таблицы \ref{tbl:time} и рисунка \ref{fig:time}, оптимизированный алгоритм оказался более эффективным.

\section{Вывод}

Из исследования видно, что неоптимизированный алгоритм Винограда уступает в эффективности стандартному, однако при его оптимизации получается получить выигрыш по времени.

\clearpage
